\chapter{Contexte et cadre du stage}
\label{chap:contexte}
Après avoir présenté l'organisme d'accueil, ce chapitre expose le cadre de mon
stage~: les objectifs et les missions qui m'ont été confiés, le cahier des
charges du projet principal, mon intégration et ma prise de poste,
l'organisation agile de l'équipe, puis l'environnement technique.

\section{Objectifs et missions}
\label{sec:objectifs}
J'ai rejoint le service Collaborative Solutions en tant qu'analyste développeur
.NET~/ Angular. La fiche de mission résumait ainsi le poste~:
\begin{quote}\itshape
Concevoir, développer, maintenir et gouverner des solutions applicatives
autour de l'écosystème Microsoft~365, et particulièrement SharePoint~Online.
En tant qu'analyste développeur, vous interviendrez sur l'ensemble du cycle de
vie des applications.
\end{quote}
Ce périmètre couvre quatre volets complémentaires. Le développement
d'applications, avec une participation aux choix d'architecture et
l'intégration des normes d'éco-conception et d'accessibilité numérique.
L'analyse et la conception, qui regroupent l'analyse des besoins métier, la
rédaction des spécifications techniques et l'amélioration continue des
solutions existantes. Les tests et la qualité, à travers les tests unitaires
et le respect des bonnes pratiques. Enfin la maintenance corrective et
évolutive, incluant la résolution d'incidents en production.

En pratique, mes missions se sont organisées autour d'un projet principal et de
plusieurs missions transverses. Le projet principal a été la refonte de
l'application interne \csbag, dont la version~2 a été découpée en trois lots
complémentaires, présentés dans la section suivante.

Les missions transverses, plus courtes mais représentatives du travail
d'exploitation en production, ont porté sur deux applications existantes. Elium
Sauvegarde est une application console .NET qui archive automatiquement les
articles de la plateforme de gestion des connaissances Elium vers
SharePoint~Online~; j'y suis intervenu pour fiabiliser un traitement qui
échouait en production. Access Report TTC~GO est un ensemble d'Azure WebJobs qui
génèrent à la demande un rapport des droits d'accès d'un espace collaboratif
SharePoint~; j'y ai corrigé une panne silencieuse d'authentification. Ces deux
missions sont détaillées au chapitre~\ref{chap:csbag}.

\section{Cahier des charges}
\label{sec:cahier-charges}
\csbag est un outil interne au service Collaborative Solutions, destiné aux
profils techniques de l'équipe. Il sert à recenser et documenter les
applications du service ainsi que leurs composants techniques~: frameworks,
comptes de service, applications SharePoint et Azure~AD, certificats,
ressources Azure, serveurs d'exploitation et bases de données. Son accès est
réservé aux comptes d'administration (comptes ADM)\footnote{Source~:
documentation interne VINCI~Energies Systèmes d'Information.}.

L'objectif de la version~2 est de faire évoluer cet outil pour en améliorer la
richesse fonctionnelle, l'ergonomie et la maintenabilité. Pour être menée de
façon progressive et livrable par étapes, la refonte a été découpée en trois
lots complémentaires, chacun formalisé par un document de spécifications
propre. Cette approche incrémentale permet de livrer de la valeur à chaque
étape tout en limitant les risques. Le détail technique de chaque lot est
présenté au chapitre~\ref{chap:csbag}.

\begin{figure}[H]
  \centering
  \begin{tikzpicture}[
    font=\footnotesize,
    lot/.style={rectangle, rounded corners, draw=vinciblue, very thick,
      fill=vinciblue!8, text width=3.1cm, minimum height=1.7cm, align=center,
      inner sep=5pt},
    num/.style={circle, fill=vinciblue, text=white, font=\bfseries\small,
      inner sep=1pt, minimum size=5.5mm},
    fleche/.style={-{Stealth[length=3mm]}, draw=vincired, very thick}
  ]
    \node[lot] (l1) {Socle\\ de la refonte};
    \node[lot, right=1.3cm of l1] (l2) {Amélioration\\ de l'outil};
    \node[lot, right=1.3cm of l2] (l3) {Enrichissement\\ fonctionnel};
    \node[num] at (l1.north) {1};
    \node[num] at (l2.north) {2};
    \node[num] at (l3.north) {3};
    \draw[fleche] (l1) -- (l2);
    \draw[fleche] (l2) -- (l3);
    \node[below=0.5cm of l2, font=\itshape, text=grisfonce, text width=11cm,
          align=center]
      {Livraison incrémentale~: chaque lot s'appuie sur le précédent
       et étend le périmètre fonctionnel.};
  \end{tikzpicture}
  \caption{Découpage de la refonte \csbag~v2 en trois lots complémentaires.}
  \label{fig:lots}
\end{figure}

\section{Intégration et prise de poste}
\label{sec:integration}
Le stage a débuté par deux journées d'accueil organisées par les ressources
humaines de VESI, communes à l'ensemble des stagiaires. Après un tour de
présentation, ces journées ont enchaîné plusieurs interventions des différentes
équipes et directions de VESI, permettant de comprendre leurs rôles et
d'identifier les interlocuteurs à solliciter. Elles se sont conclues par un
échange avec deux directeurs, l'occasion de recueillir des conseils sur le
parcours professionnel.

À partir du troisième jour, j'ai rejoint mon équipe sur le site du Mans, où
j'ai été accueilli par mon chef de projet, Maxime, et le team lead de l'équipe,
Alexis. Après la découverte des locaux, un premier point avec mon chef de
projet a permis de fixer la feuille de route du stage. La première semaine a
ensuite été consacrée à la préparation de l'environnement de travail et à la
validation des formations et passeports obligatoires, avant le lancement du
premier lot.

\begin{table}[H]
  \centering
  \caption{Étapes de la prise de poste.}
  \label{tab:integration}
  \begin{tabularx}{\textwidth}{>{\bfseries}p{3.4cm} X}
    \toprule
    \normalfont\bfseries Étape & \bfseries Contenu \\
    \midrule
    Journées d'accueil (J1--J2) & Présentation des équipes et directions de
      VESI, règles internes, échange avec deux directeurs \\
    Prise de poste (J3) & Accueil dans l'équipe, découverte des locaux,
      cadrage de la feuille de route avec le chef de projet \\
    Semaine 1 & Installation de l'environnement de travail~; formations et
      passeports obligatoires (sécurité des systèmes d'information, règles
      d'usage de l'IA~: seul GitHub~Copilot autorisé) \\
    Lancement du lot 1 & Présentation du projet \csbag et démarrage des
      développements \\
    \bottomrule
  \end{tabularx}
\end{table}

\section{Organisation du travail et méthode agile}
\label{sec:organisation}
Au sein de la direction Custom Applications, le service Collaborative Solutions
comprend l'équipe Group Collaborative Solutions, elle-même organisée en deux
squads~: le squad MyVIEW et le squad M365. J'ai été intégré au squad M365.

\begin{figure}[H]
  \centering
  % >>> Remplacer la ligne ci-dessous par votre export draw.io : <
   \includegraphics[width=\textwidth]{org_squads.png}
  \caption{Positionnement du squad M365 au sein de la direction Custom
    Applications.}
  \label{fig:org-squads}
\end{figure}

L'équipe fonctionne selon une approche agile, rythmée par plusieurs rituels au
cours de la semaine. Chaque jour, un daily réunit le squad pour faire le point
sur le travail de la veille, celui du jour, et remonter d'éventuels points
bloquants. Le lundi, un point d'équipe rassemble l'ensemble de l'équipe Group
Collaborative Solutions afin de coordonner les travaux entre les deux squads.
Le mardi, un point de squad se concentre sur l'avancement propre au squad M365.
Enfin, le vendredi, un point technique avec le référent technique, Nathan,
permet de valider les choix techniques et de discuter en profondeur des sujets
de conception avant leur mise en œuvre. Le suivi des tâches s'effectue tout au
long de la semaine dans Wrike, où le travail est formalisé et priorisé, ce qui
assure la traçabilité des développements et une visibilité partagée sur
l'avancement.

\begin{table}[H]
  \centering
  \caption{Rituels agiles suivis durant le stage.}
  \label{tab:rituels}
  \begin{tabularx}{\textwidth}{>{\bfseries}p{3.2cm} p{2.6cm} X}
    \toprule
    \normalfont\bfseries Rituel & \bfseries Fréquence & \bfseries Objectif \\
    \midrule
    Daily          & Quotidienne & Point sur la veille et la journée à venir,
                                   remontée des points bloquants \\
    Point d'équipe & Hebdo (lundi) & Coordination à l'échelle de l'équipe
                                     Group Collaborative Solutions \\
    Point de squad & Hebdo (mardi) & Suivi de l'avancement au sein du
                                     squad M365 \\
    Point technique & Hebdo (vendredi) & Validation des choix techniques avec
                                         le référent technique \\
    \bottomrule
  \end{tabularx}
\end{table}
\section{Environnement technique}
\label{sec:environnement-technique}
Les développements s'appuient sur une pile technique full-stack cohérente avec
celle du service, complétée par des outils de gestion de code, de suivi de
projet et d'assistance au développement.

\begin{table}[H]
  \centering
  \caption{Environnement technique du stage.}
  \label{tab:stack}
  \begin{tabularx}{\textwidth}{>{\bfseries}p{3.8cm} X}
    \toprule
    \normalfont\bfseries Catégorie & \bfseries Technologies et outils \\
    \midrule
    Back-end          & C\#, ASP.NET~Core (.NET~Core / .NET~Framework) \\
    Front-end         & Angular \\
    Base de données   & Microsoft SQL~Server \\
    Cloud             & Microsoft Azure (services cloud, Azure~AD) \\
    Gestion de code   & GitLab (instance auto-hébergée en interne) \\
    Gestion de projet & Wrike \\
    Assistance IA     & GitHub~Copilot (seul outil d'IA autorisé) \\
    Rédaction du rapport & Latex via la plateforme Overleaf~\cite{overleaf} \\
    \bottomrule
  \end{tabularx}
\end{table}
Ce socle est représentatif des projets du service~: un back-end .NET exposant
des services consommés par un front-end Angular, une persistance sur SQL~Server
et une intégration à l'écosystème Microsoft~365 et Azure. Il constitue le cadre
technique dans lequel s'inscrit la refonte de \csbag, présentée au
chapitre~\ref{chap:csbag}.